\section{Pipeline de generación de conjuntos de datos}

\subsection{Generación de Instruction Data}

El expositor presentó un pipeline típico de generación de instruction data:

\begin{enumerate}
    \item \textbf{Seed Data}: el punto de partida puede ser texto original, pares de preguntas y respuestas (QA) ya existentes, o consultas reales de usuarios
    \item \textbf{Generación de instrucciones}: si solo se cuenta con texto original (es decir, la respuesta), se puede usar \textbf{back translation} (traducción inversa) para que el LLM genere la pregunta correspondiente
    \item \textbf{Generación de respuestas}: se usa un LLM para generar la respuesta a partir de la instrucción
    \item \textbf{Evaluación y filtrado}:
    \begin{itemize}
        \item Se evalúa la calidad con reglas heurísticas o con LLM-as-Judge
        \item Deduplicación de los datos (deduplication)
        \item Descontaminación (decontamination): garantizar que no se incluyan muestras del conjunto de prueba
        \item Filtrado por palabras clave: retirar contenido no deseado
    \end{itemize}
\end{enumerate}

\subsubsection{Tratamiento especial de los datos de Instruction Following}

Instruction following es un tipo de datos centrado en el \textbf{cumplimiento de restricciones}. Por ejemplo: «responde en inglés y todo en minúsculas». Al generar este tipo de datos, se puede anexar la condición de restricción al final del prompt para que el LLM genere la respuesta, y luego usar \textbf{pruebas programáticas} para verificar si se cumple la restricción (como detección de idioma o verificación de mayúsculas y minúsculas), descartando las muestras que no cumplan los requisitos.

\subsection{Generación de Preference Data: el paradigma Ultra Feedback}

La generación de Preference data sigue una estrategia distinta:

\begin{enumerate}
    \item Dado un prompt, se consulta a \textbf{varios LLM distintos} (en lugar de un solo modelo)
    \item Se usa un Judge LLM para calificar cada respuesta
    \item Se toma la respuesta con la calificación más alta como \textbf{chosen} y la de calificación más baja como \textbf{rejected}
    \item Deduplicación, retiro de respuestas demasiado cortas y otros procesamientos posteriores
\end{enumerate}

\begin{knowledgebox}{Decontamination (descontaminación)}
La descontaminación es un paso fundamental en el pipeline de generación de datos. Por ejemplo, IFEval es un conjunto de evaluación de instruction following ampliamente usado. Al generar datos de entrenamiento, es indispensable garantizar que las muestras generadas no tengan una alta similitud con el conjunto de prueba; de lo contrario, la puntuación de evaluación se inflará de forma artificial sin que la capacidad real mejore.
\end{knowledgebox}

\subsection{Caso: el conjunto de datos Open Perfect Blend}

El expositor mostró el conjunto de datos Open Perfect Blend que elaboró, un conjunto de datos mixto orientado al fine-tuning general. Su proporción aproximada por categoría es:

\begin{itemize}
    \item \textbf{Math} (matemáticas): proporción mayor
    \item \textbf{Code} (código): proporción mayor
    \item \textbf{Chat} (diálogo general): proporción media
    \item \textbf{Instruction Following}: proporción menor
\end{itemize}

El diseño de esta proporción de mezcla (data mixture) determina directamente el desempeño del modelo en las distintas dimensiones de capacidad.

\subsection{Chat Template}

Una vez terminada la preparación de los datos, el último paso es aplicar el \textbf{Chat Template}. Esta es la estructura clave que permite al modelo entender «quién está hablando»:

\begin{lstlisting}[caption={Ejemplo de Chat Template (formato ChatML)}]
<|im_start|>system
You are a helpful assistant.
<|im_end|>
<|im_start|>user
What is the capital of France?
<|im_end|>
<|im_start|>assistant
The capital of France is Paris.
<|im_end|>
\end{lstlisting}

Los tokens especiales \texttt{<|im\_start|>} y \texttt{<|im\_end|>} marcan el inicio y el fin de cada mensaje, y las etiquetas de rol (system / user / assistant) precisan la identidad de quien habla. Las conversaciones de varios turnos se logran repitiendo el par user-assistant.

\begin{importantbox}{El Chat Template es el objetivo central de aprendizaje del SFT}
La estructura del Chat Template es uno de los contenidos centrales que el modelo debe aprender en la etapa de SFT. Es precisamente este formato de conversación estructurado el que hace que el base model evolucione de «solo saber continuar texto» a «un assistant capaz de sostener conversaciones de varios turnos». Distintas familias de modelos (Llama, Mistral, Qwen, etc.) usan chat templates diferentes, pero el principio es el mismo.
\end{importantbox}

\subsection{Resumen de la sección}

La generación de conjuntos de datos es un pipeline sistemático que incluye la preparación de los datos semilla, la generación de instrucciones y respuestas, la evaluación y el filtrado, y la deduplicación y descontaminación, entre otros pasos. Los datos de preferencia (Preference data) se obtienen mediante muestreo con varios modelos más la calificación de un Judge. Finalmente, los datos deben formatearse con el Chat Template para usarse en el entrenamiento de SFT.

\newpage
